\documentclass[11pt]{article}
\usepackage[T1]{fontenc}
\usepackage{lmodern,microtype,amsmath,amssymb,graphicx,geometry,booktabs,array,xcolor,caption,hyperref,float}
\geometry{letterpaper,margin=0.78in,headheight=14pt}
\definecolor{ink}{HTML}{252C30}\definecolor{accent}{HTML}{A84B32}
\hypersetup{colorlinks=true,urlcolor=accent,linkcolor=ink,pdftitle={Inferring a recording location from aircraft sound},pdfauthor={James N. Jean IV}}
\captionsetup{font=small,labelfont=bf,labelsep=period}
\setlength{\parindent}{0pt}\setlength{\parskip}{6pt}\setlength{\emergencystretch}{2em}
\newcommand{\target}{ A-27}
\newcommand{\fig}[3]{%
\begin{figure}[H]\centering
\IfFileExists{figures/#1.pdf}{\includegraphics[width=\linewidth,height=#2,keepaspectratio]{figures/#1.pdf}}{%
\IfFileExists{figures/#1.png}{\includegraphics[width=\linewidth,height=#2,keepaspectratio]{figures/#1.png}}{%
\fbox{\parbox[c][#2][c]{0.9\linewidth}{\centering Figure pending}}}}
\caption{#3}\end{figure}}

\newcommand{\nextsection}[1]{\clearpage\section{#1}}
\begin{document}
\begin{center}
{\LARGE\bfseries Inferring a recording location\\[3pt] from aircraft sound}\par
\vspace{7pt}
{\large Doppler measurements, flight-path comparison, and observer geometry}\par
\vspace{7pt}
{\large James N. Jean IV}\par
\vspace{7pt}
{\small Public research manuscript: anonymized measured evidence}
\end{center}
\begin{abstract}
A passing aircraft leaves a changing harmonic pattern that can constrain a recording location. We analyze two clips of one flyover, preserving their 15.160\,s gap and comparing measured Doppler structure with public aircraft trajectories. Capture timestamps establish the recording window. A search of the full archive at that window used neither an aircraft-identity filter nor an observer-location prior. Among 358 refined candidates ranked under each source model, the supplied aircraft ranks first by curve RMS with constant pitch and bounded pitch drift. Its corresponding location estimates are 56.5\,m and 79.1\,m from the recorder's approximate reference; spectral-contrast alternatives differ. These results support corrected-time model agreement, not a retrospective blind success: aircraft identity and reference location had already been disclosed, the reference accuracy is unknown, and nearby alternative solutions remain. 
\end{abstract}
\fig{measured_curve}{2.55in}{Measured common harmonic spacing on the preserved elapsed-time axis. The blank interval is unrecorded, not observed silence. The spacing convention is not a verified engine fundamental.}
\section{Question and experimental scope}
The forward problem predicts the received pitch from an aircraft trajectory and a listener position. The inverse problem searches for trajectories and positions consistent with the measured curve. This manuscript follows that progression from audio to geometry, candidate comparison and location. The corrected search was conducted after disclosures, but its code did not receive the disclosed aircraft identity or recording coordinates as search inputs. That procedural separation does not make the experiment newly blind.

\nextsection{Extracting a curve without inventing the gap}
The clips contain one flyover, not two independent acoustic events. The first lasts 69.840\,s; the second lasts 75.465\,s and begins 85.000\,s after the first start. The resulting gap is 15.160\,s and the complete elapsed span is 160.465\,s. Approximately one-second gap perturbations were tested as timestamp-quantization sensitivity, not as an independently established accuracy bound. The absent interval cannot be used as evidence of silence.

\fig{measurement_logic}{3.35in}{Measurement logic: the left-hand $4f$ and $8f$ curves illustrate harmonic assignments reconstructed from the measured common spacing; they are not independently extracted band traces. The arithmetic example is schematic. The right-hand 62 points are measured common-spacing observations. Raw audio and original spectral pixels are withheld from the public edition.}

Sixty-two reviewed time blocks describe the shared frequency trend. If a ridge is assigned harmonic index $h$, its frequency $f_h$ contributes to the common-spacing measurement $y=f_h/h$. This makes different ridges comparable, but it does not make their errors independent: harmonics observed in the same audio share source modulation, interference and labeling uncertainty.

A second objective uses the spectrogram directly rather than only the traced blocks. Seventy non-overlapping 2.048\,s frames provide frequency-local contrast along predicted harmonic curves. These are two ways to score the same recording, not independent validation datasets. Persistent tones, overlapping sources and a chosen harmonic labeling can all affect the result. Accordingly, neither a good curve fit nor a bright predicted ridge is assigned an identity probability.

\subsection*{Could the sound be recorded indoors?}
Low-frequency aircraft sound can reach indoors. The FAA-funded PARTNER study reports greater outdoor-to-indoor transmission at low than high frequencies~\cite{partner}; FAA guidance addresses sound insulation of buildings exposed to aircraft noise~\cite{insulation}. Buildings, windows, room resonances, microphones and noise affect recoverable harmonics. Audible rumble alone is insufficient: the method needs a stable, attributable Doppler ridge. The PARTNER work largely concerns airport/runway noise, not indoor geolocation, and indoor Doppler recovery has not been validated in this study. Paired, synchronized indoor/outdoor recordings should test ridge preservation before extending the method.

\nextsection{From frequency change to three-dimensional geometry}
Let $\mathbf{s}(t_e)$ be aircraft position at emission time $t_e$ and $\mathbf{x}$ the stationary listener position. The slant distance is
\begin{equation}
 R(t_e)=\|\mathbf{x}-\mathbf{s}(t_e)\|,\qquad
 t_r=t_e+\frac{R(t_e)}{c},
\end{equation}
where $t_r$ is reception time and $c$ is sound speed. In still air, the received frequency is
\begin{equation}
 f_r(t_r)=\frac{f_s(t_e)}{1+\dot R(t_e)/c}.
 \label{eq:doppler}
\end{equation}
Approach gives $\dot R<0$ and raises the received frequency; recession lowers it. Crucially, the aircraft position is evaluated when the sound was emitted, not when the microphone received it.
\fig{geometry}{3.15in}{Two conceptual listener candidates at different ground elevations, with the same aircraft altitude. Not to scale; the 1.5 m observer height is exaggerated. MSL denotes a common approximate orthometric datum. Actual ranges use Earth-centered, Earth-fixed coordinates; curvature, geoid variation and approximate pressure-height conversion prevent treating flat altitude subtraction as exact range geometry.}
The environmental model uses a uniform wind $\mathbf{w}$. With $\mathbf{q}=\mathbf{x}-\mathbf{s}(t_e)$, delay $\tau$ solves $\|\mathbf{q}-\mathbf{w}\tau\|=c\tau$. Defining the air-relative ray direction $\mathbf{n}=(\mathbf{q}-\mathbf{w}\tau)/(c\tau)$ gives
\begin{equation}
 D=\frac{f_r}{f_s}=\left[1-\frac{\mathbf{n}\cdot\mathbf{v}_{\rm ground}}{c+\mathbf{n}\cdot\mathbf{w}}\right]^{-1}.
\end{equation}
Delay and frequency thus use one consistent moving-air model. Uniform wind is an approximation: spatial gradients and atmospheric refraction are not modeled.

\nextsection{A trajectory search at the capture time}
Capture metadata establish the recording start; the second clip starts 85.000\,s later. Explicit timezone tags agree across the original video streams, although the device clock was not externally synchronized. A $\pm2$\,s clock-offset range around the capture timestamp is a modeling choice, not a calibrated clock-error interval.

\fig{flow_counts}{3.2in}{Corrected-time search scope. The 361-aircraft refinement set is the union of the declared top 200 candidates for each of two source models and two coarse objectives. It is not an identity or type filter. Final ranks are among refined candidates, not an exhaustive proof over all possible locations.}
The source archive contains 75,222 files and 101,298,620 raw points. Within the capture window, 6,743 aircraft have relevant airborne reports; 3,880 produce covered coarse modes. The flag-aware scan evaluates 9,331,875 location--time pairs. A declared union rule selects 361 aircraft for environmental refinement, which evaluates another 1,859,836 pairs. Of these, 359 have retained valid modes in at least one model, and 358 are ranked in each source-model/objective combination.

The source-window screen extends from 600 s before capture start to 800 s after it, retaining delayed-source coverage rather than testing only receive-time airborne activity. Coarse ridge seeding uses 17 of the 62 reviewed blocks; final ridge scoring uses all 62.

Reported ground states and gaps exceeding 60\,s split segments; stale/leg markers are not treated as invalid positions. Geometric altitudes are not pressure-corrected. These coverage rules matter because rejecting short or flagged tracks without understanding the flags can remove plausible candidates.

\nextsection{Aircraft types: comparison, not acoustic classification}
The refinement set spans 108 archive-reported type categories among 361 aircraft. It includes 33 P28A, 31 C172, 17 B738, 13 C208, 12 PC12, 11 B763 and 10 AT76 aircraft, plus 10 unknown types. No type exclusion is applied to the main ranking. Metadata describe the archive entries; they do not independently classify the recorded engine sound.
\fig{type_comparison}{3.25in}{Type composition and candidate comparisons in the fresh refinement set. The target is not selected by filtering to its type. Unknown types remain in the search. Categories are archive-reported, not independently registry-verified.}
Periodic sound is compatible with several mechanisms. Blade-passage frequency is $f_{\rm BPF}=B\,\mathrm{RPM}/60$; a simple four-stroke ignition cadence is $f_{\rm ign}=N_c\,\mathrm{RPM}/120$. A two-blade propeller and four-cylinder four-stroke engine both produce 80\,Hz at 2400\,RPM. This overlap prevents the observed harmonics from uniquely identifying engine configuration or aircraft type.

Post-freeze gates are sensitivity tests, not main filtering. A conventional-jet-code exclusion leaves 257/258 candidates for constant/drifting source models; a broad computational speed/range screen leaves 258/260. An illustrative 50--90\,m/s, 1.0--5.3\,km CPA screen leaves 204/213. A tighter 60--70\,m/s, 1.73--2.84\,km screen leaves 55/52 and rejects the target. The target is already first by RMS without any of these gates. A narrower filter is not necessarily a better one.

\nextsection{Ranking the competing trajectories}
For model parameters $\theta$, measured common spacing $y_i$ and prediction $\widehat y_i$, the ridge objective is
\begin{equation}
E(\theta)=\sqrt{\frac{1}{62}\sum_{i=1}^{62}[\widehat y(t_i;\theta)-y_i]^2}.
\end{equation}
The source is either constant in emitted spacing or allowed a bounded change,
\begin{equation}
f_s(t_e)=f_{\rm mid}\left[1+\delta\frac{t_e-t_{\rm mid}}{T_e}\right],\qquad |\delta|\leq0.03.
\end{equation}
Here $T_e$ is the fitted emission span: $\delta$ is the end-to-end fractional change relative to midpoint spacing. The free source spacing prevents a harmonic label from serving as an assumed aircraft identity.
\fig{score_scatter}{3.0in}{Refined candidate scores under constant and drifting source models. The target's strongest ridge fits do not imply it wins every objective. Each plotted aircraft pairs its best RMS with its best spectral contrast; the extrema can come from different retained geometry modes.}
\begin{center}\small
\begin{tabular}{@{}lrrr@{}}\toprule
Source model & Target RMS & RMS rank & Profiled pixel rank\\\midrule
Constant & 0.109325\,Hz & 1/358 & 2/358\\
Up to 3\% drift & 0.095596\,Hz & 1/358 & 5/358\\\bottomrule
\end{tabular}
\end{center}
Pixel ranks use equal nine-point $\pm1\%$ spacing profiling. They compare retained, ridge-screened geometry modes rather than an independent exhaustive pixel optimization. Neither a residual nor a contrast score is a calibrated identity probability.

\nextsection{From coarse geometry to an environmental fit}
The search uses finite trajectory-centered spatial domains rather than an observer-location prior. Coarse seeding uses a 1\,km grid over a $\pm12$\,km domain and still air at 343\,m/s. Selected candidates receive 250\,m and 100\,m refinement with terrain and atmospheric inputs. Refinement selection combines coarse ridge and pixel objectives; it does not simply follow the eventual winning aircraft.
\fig{curve_fit}{3.35in}{Measured common spacing and corrected-capture model predictions, with residuals. The two source models are compared on the same observed blocks.}
Terrain is sampled from ETOPO2022 with an EGM2008 geoid companion. Native spacing is approximately 1.85\,km north--south, so 100\,m search steps do not create 100\,m terrain information. Negative relief is floored at zero and listener height is ground plus 1.5\,m. GFS fields at 0.25-degree spacing supply uniform effective wind, temperature-dependent sound speed and pressure-height corrections. These are coarse modeled conditions, not a surveyed local atmosphere. Reported geometric heights bypass the pressure-altitude conversion.

The coarse pixel screen uses 18 frames; the final comparison uses 70 non-overlapping 2.048\,s frames. Whitened power near predicted fourth/eighth harmonics is contrasted with nearby control lines, balanced across clips/channels and trimmed. This direct spectral evidence complements the 62 reviewed blocks but reuses the same audio; it is not independent validation.

\nextsection{A location estimate from the measured sound}
The corrected inverse fits place the listener near the approximate recording reference. The reference is not used to move the frozen corrected estimates. Its agreement is evaluated only after the corrected ranking is frozen. Because the reference and aircraft identity were already known to the investigators, this is a post-disclosure comparison rather than a fresh held-out or blind test.

\fig{observer_geometry_schematic}{2.5in}{Non-geographic geometry illustration. The actual path, estimated coordinates and approximate recording reference are withheld publicly; scalar separation errors are reported in Section~10. This diagram is schematic, not a map of the event.}
The public edition preserves model definitions and scalar results but omits recording coordinates, flight identifiers, geographic orientation and actual trajectory maps. The corrected location comparison can be understood without disclosing the recording area.

At each RMS-selected mode's continuous closest approach, the model gives:
\begin{center}\small
\begin{tabular}{@{}lrr@{}}\toprule
Quantity (m) & Constant source & Bounded drift\\\midrule
Aircraft altitude, approximate MSL & 1392.04 & 1392.25\\
Candidate ground elevation, approximate MSL & 352.39 & 352.78\\
Observer height above ground & 1.50 & 1.50\\
Local-tangent vertical separation & 1038.05 & 1037.85\\
Horizontal separation & 1140.27 & 1237.55\\
Slant range & 1542.00 & 1615.13\\\bottomrule
\end{tabular}
\end{center}
These are model-derived heights and ranges, not surveyed elevations. For example, the constant-source orthometric altitude difference is 1038.15 m, slightly different from the local-tangent 1038.05 m because of curvature and geoid variation.

Continuous CPA minimization over the recorded emission interval gives slant ranges of 1.542\,km for the constant-source RMS fit and 1.615\,km for the drifting-source RMS fit. Corresponding horizontal distances are 1.140 and 1.238\,km; local-tangent vertical clearance is about 1.038\,km. These are distinct geometric quantities, not interchangeable versions of altitude. The Earth-fixed speed magnitude is approximately 56.68\,m/s and includes vertical motion. These values are conditional model outputs, not direct range measurements.

\nextsection{Alternative basins and forward consistency}
After freezing the corrected ranking, a target-only profile explores a $\pm12$\,km domain at 250\,m spacing with the same $\pm2$\,s clock bounds. This is a diagnostic of the selected trajectory, not another all-aircraft identification test and not a calibrated location posterior.
\fig{model_sensitivity}{3.3in}{Minimum RMS increase by spatial-distance bin in the post-freeze target-only profile. The coarse profile exposes alternative basins; it does not establish a confidence radius or guarantee that all possible solutions have been found. Its constant-pitch grid minimum need not equal the finer frozen-search minimum.}
More than 1\,km from the profile minimum, the best alternative is worse by only 0.04189\,Hz under constant pitch and 0.02687\,Hz under drifting pitch. A secondary drift basin about 2.5\,km away therefore remains competitive on the scale of an uncalibrated residual. Small error against one approximate reference should not be converted into a general accuracy claim.

A separate forward check fixes the listener at the disclosed reference and profiles nuisance parameters. RMS is 0.111969\,Hz for constant pitch and 0.099394\,Hz for drifting pitch, close to the corresponding free-location fits. The constant-source clock hits its $+2$\,s bound. This supports consistency with the reference under the model; it is not an independent prediction, since the location is fixed using disclosed information. Source uncertainty, model weather, clock calibration and harmonic correlations remain relevant.

\nextsection{Corrected location agreement and its limits}
\fig{corrected_location_agreement}{3.35in}{WGS84 geodesic separations between frozen corrected-search estimates and the supplied approximate recording point. These are post-disclosure agreement measurements, not held-out accuracy guarantees. Each source model contributes its RMS-selected and spectral-contrast-selected estimate.}
\begin{center}\small
\begin{tabular}{@{}llr@{}}\toprule
Source model & Location selected by & Separation from reference\\\midrule
Constant & Curve RMS & 56.5\,m\\
Up to 3\% drift & Curve RMS & 79.1\,m\\
Constant & Spectral contrast & 167.6\,m\\
Up to 3\% drift & Spectral contrast & 158.6\,m\\\bottomrule
\end{tabular}
\end{center}
The two RMS-selected estimates agree with the approximate reference within 80\,m. That is meaningful agreement for this corrected-time reconstruction, but the point's accuracy has not been independently established. If it is within $\epsilon$ of the true recorder, each separation differs from the true location error by at most $\epsilon$ by the triangle inequality. The unknown $\epsilon$ is not a fitted confidence bound.

The reference was known before this corrected analysis, even though its coordinates were omitted from the search inputs. This is consequently a post-disclosure comparison, not a blind validation. It supports the narrower claim that a fresh, identity-unfiltered corrected-time search yields a best-RMS trajectory and conditional locations consistent with the supplied reference. The alternative pixel leaders and spatial basins prevent a unique-identity or general meter-level accuracy claim.

\nextsection{Conclusions and reproducibility}
With the capture time corrected, the measured Doppler curve produces a strong model-conditional result: the supplied aircraft ranks first by RMS in both source models among 358 refined candidates, and its RMS-selected listener estimates are 56.5\,m and 79.1\,m from the approximate recording point. The analysis uses sound to constrain observer geometry rather than inserting that point into the search. 

The limits are equally specific. The corrected analysis was post-disclosure, the reference is approximate, refinement is finite and selected, pixel ranking is restricted to retained geometry modes, and alternative spatial basins remain. None of these results constitutes a calibrated identity probability or a validated 80\,m accuracy specification. 

An independent numerical reproduction across all 359 covered aircraft agreed within $8.9\times10^{-16}$ Hz for RMS and $2.4\times10^{-13}$ dB for pixel contrast. This verifies numerical reproducibility, not physical accuracy.

A stronger next experiment should verify capture timestamps independently; audit aircraft identifiers and trajectory coverage; use separated flyovers or synchronized receivers; and calibrate source-pitch and atmospheric assumptions. Methods, candidate rules and a location prediction should be frozen before a new reference position is revealed. Such a prospective test can measure error and false matches without conflating post-disclosure reconstruction with blind performance.

\subsection*{Reproducibility and publication scope}

Public materials retain anonymized measured curves, candidate scores, declared rules, scalar comparison errors and synthetic physics checks. They omit raw audio, original spectral pixels, flight identifiers, capture dates and actual geographic tracks or reference coordinates. The standalone public \LaTeX{} source contains no inactive private branches. Exact-input geographic reproduction remains private.

\begin{thebibliography}{9}\small
\bibitem{archive} ADSB.lol contributors. Public aircraft trajectory archive and associated metadata. \url{https://www.adsb.lol/}. Coverage and metadata are imperfect; archive membership is not proof of acoustic audibility.
\bibitem{etopo} NOAA National Centers for Environmental Information. ETOPO 2022 global relief model and EGM2008 geoid companion. \url{https://doi.org/10.25921/fd45-gt74}.
\bibitem{gfs} NOAA/NCEP. Global Forecast System atmospheric model fields. \url{https://www.nco.ncep.noaa.gov/pmb/products/gfs/}. Forecast/model fields were used, not local station observations.
\bibitem{demo} Original audio-geolocation demonstration. \url{https://www.youtube.com/watch?v=UytYrhYLOKI}. The present work adapts its acoustic/trajectory reasoning; it is not an exact replication. Independent two-event timing, an observed intervening-silence test and visual corroboration were not available or were deliberately excluded here.
\bibitem{partner} Hodgdon, K. K., Atchley, A. A., and Bernhard, R. J. \textit{Low Frequency Noise Study}. PARTNER-COE-2007-001, April 2007, p.~4. FAA-funded study. \url{https://ascent.aero/documents/2020/01/low-frequency-noise-study.pdf}.
\bibitem{insulation} Federal Aviation Administration. \textit{Guidelines for Sound Insulation of Structures Exposed to Aircraft Noise}. Advisory Circular 150/5000-9B, 8 June 2022. \url{https://www.faa.gov/documentLibrary/media/Advisory_Circular/150-5000-9B-Sound-Insulation-2022.pdf}.
\end{thebibliography}
\end{document}
